\begin{pastexam}[label={px:ust-113-ee-ds-15}]{臺聯大電機類 113 學年度 第 15 題}
（複選；每題 5 分，每個選項單獨計分，答錯倒扣 1 分，倒扣至本大題 0 分為止）

In the context of graph theory and finding optimal paths or trees, which statements accurately
describe features of Dijkstra's, Prim's, and Kruskal's algorithms? (Select all that apply)
\begin{enumerate}[label=(\Alph*)]
\item Dijkstra's Algorithm operates on graphs with non-negative edge weights only.
\item Kruskal's Algorithm guarantees finding a minimum spanning tree for both directed and
undirected graphs.
\item Prim's Algorithm selects the next node based on the lowest edge weight.
\item Dijkstra's Algorithm guarantees finding the longest path from a single source to all other
vertices in a graph.
\item Prim's Algorithm guarantees the creation of a minimum spanning tree starting from an
arbitrary vertex.
\end{enumerate}
\end{pastexam}

\begin{pastsoln}
答案：(A)、(C)、(E)。

\begin{itemize}
\item (A) Dijkstra 的正確性證明用到「已經確定的頂點，距離不會再被更短的路徑改小」，這需要邊的權重都
$\ge 0$；有負邊時可能算錯（資結書第~25.2~節，有負邊要用 Bellman-Ford）。 對。
\item (B) Kruskal 處理的是無向圖。 有向圖的「最小生成樹」是最小樹形圖（從根出發到每個頂點都有有向路徑），
Kruskal 不保證找得到。 例如三個頂點、四條有向邊：$a \to r$ 權重 $1$、$a \to b$ 權重 $2$、$r \to a$ 權重 $5$、
$r \to b$ 權重 $6$，以 $r$ 為根。 Kruskal 不看方向，先選 $a \to r$、再選 $a \to b$（不成環），就有 $2$ 條邊、
總權重 $3$；但 $r$ 沒有任何出去的邊，從 $r$ 到不了 $a$、$b$，不是以 $r$ 為根的樹形圖。 正確答案是
$r \to a$、$a \to b$，總權重 $7$（要用 Chu–Liu/Edmonds 演算法）。 錯。
\item (C) Prim 每一步從「一端在樹上、一端不在樹上」的邊裡選權重最小的，把那個頂點加進樹（資結書第~24.3~節）。
選項的說法簡略，但描述的就是這件事。 對。
\item (D) Dijkstra 求的是\textbf{最短}路徑。 一般圖的最長簡單路徑是 NP 困難的問題，Dijkstra 把 $\min$ 換成
$\max$ 也不對。 錯。
\item (E) 割的性質保證每一步加入的邊都屬於某棵最小生成樹，和起點無關（資結書定理~24.7）；連通圖從任何
頂點開始都會得到最小生成樹。 對。
\end{itemize}
\end{pastsoln}
